\section{总结与讨论}
\label{sec:conclusion}

本文基于统一的数学框架系统研究了赤道波动力学问题，从三维层结大气的基本方程出发，逐步建立了涵盖干波理论、湿过程修正以及微扰分析的完整理论体系。
这一框架不仅在形式上统一了 Kelvin 波、Rossby 波、惯性重力波及混合 Rossby--重力波等传统分类，也为理解观测中对流耦合赤道波（CCEWs）的复杂行为提供了坚实的理论基础。

本文首先从旋转层结大气的原始方程组出发，在赤道 $\rossby$ 平面近似、静力平衡及线性化假设下，推导了描述赤道波动的基本方程组。
科氏参数随纬度的线性变化 $\coriolis = \rossby y$ 构成了赤道波导的动力学基础，使得波动能量被有效束缚在赤道附近的狭窄纬带内。

该框架的核心贡献在于建立了垂直与水平结构的严格分离。
通过 Sturm-Liouville 理论，垂直方向的本征方程定义了一组正交完备的本征函数 $\{\vertfunc_\vertmode(z)\}$，对应的本征值即为等效深度 $\eqdepth_\vertmode$。
等效深度作为连接三维垂直结构与二维浅水动力学的核心参数，其物理意义在于：每一个特定的垂直模态对应一个确定的等效深度，而等效深度又通过波速 $\ceq = \sqrt{\gravityconst\eqdepth}$ 决定该垂直模态在水平方向上的传播特性。
在刚性边界条件下，等效深度形成离散谱；若采用辐射边界条件，则等效深度构成连续谱。

固定垂直模态后，水平结构方程通过引入梯算子 $\aop$ 与 $\adop$ 被重新表述为算子代数形式。
这种处理方式类似于量子力学中谐振子问题的代数化求解，清晰地揭示了解空间的数学结构。
经向速度的本征函数由厄米函数 $\hfunc{\modenum}(y)$ 给出，其物理图像是：科氏参数的纬度变化在动力学上等效于二次型势阱，波动能量被束缚在赤道这一势阱中心附近。
因此，赤道波动方程组的完整解空间具有张量积结构 $\mathcal{H} = \mathcal{H}_z \otimes \mathcal{H}_{xy}$，任一本征态可由三个量子数 $(\vertmode, \modenum, \wavenum)$ 唯一标识：垂直量子数 $\vertmode$ 决定等效深度与垂直节点结构，经向量子数 $\modenum$ 决定经向节点数与波动类型，纬向波数 $\wavenum$ 决定纬向传播方向与空间周期。

在上述框架内，本文系统分析了干大气条件下赤道波动的本征模态。
垂直方向上，假设 Brunt-Vaisala 频率为常数，本征方程可显式求解，其解表现为振幅随高度指数增长、相位随高度振荡的内波结构。
振幅的增长源于背景密度的指数衰减，能量密度实际保持守恒。
当等效深度超过临界值时，垂直结构从振荡型过渡为指数衰减型，对应能量被限制在低层大气的俘获模态。

水平方向上，对每一组 $(\wavenum, \modenum)$，Matsuno 色散关系给出三个本征频率，分别对应两支惯性重力波与一支赤道 Rossby 波。
惯性重力波在高频极限下满足 $\frequency^2 \approx \wavenum^2 + 2\modenum + 1$，其恢复力主要来自重力与科氏力的共同作用。
赤道 Rossby 波在低频极限下满足 $\frequency \approx -\wavenum/(\wavenum^2 + 2\modenum + 1)$，始终向西传播，其恢复机制源于赤道 $\rossby$ 效应。
当 $\modenum = 0$ 时，混合 Rossby--重力波在重力波与 Rossby 波之间起连续过渡作用。
赤道 Kelvin 波作为特殊模态需单独处理，其特征是经向速度恒为零、色散关系线性无频散、仅允许向东传播。

干大气理论成功解释了平流层中观测到的 Kelvin 波、混合 Rossby--重力波及惯性重力波等基本模态。
早期观测研究在平流层识别出的波动，其传播特性与频率特征均与理论预测高度吻合，为 Matsuno 理论提供了直接的观测验证。

干大气理论虽然成功描述了平流层波动，但无法解释对流层中观测到的对流耦合赤道波（CCEWs）的慢传播特性。
观测表明，CCEWs 的相速度约为干波预测值的 $1/2$ 至 $1/3$，对应的等效深度仅为 $12$--$50$ m，远小于干大气第一斜压模态的典型值（约 $200$ m）。
这一显著偏差表明，非绝热加热 $\nonadheat$ 在热带对流层波动中起关键作用。

本文通过 Wave-CISK 参数化将湿过程效应纳入理论框架。
该参数化假设非绝热加热率与垂直上升运动强度成正比，其物理基础是热带大气中低层辐合触发深对流从而释放潜热的机制。
在数学上，这一参数化等价于将 Brunt-Vaisala 频率的平方替换为有效静力稳定度 $\bvfreq_{\mathrm{eff}}^2 = (1-\moistparam)\bvfreq^2$，从而使有效等效深度减小为 $(1-\moistparam)\eqdepth$，波速相应降低为 $\sqrt{1-\moistparam}\cdot\ceq$。
色散关系的形式保持不变，但所有模态的频率按统一因子缩放。
当 $\moistparam \approx 0.75$ 时，波速约减半，与观测结果定性一致。

这一结果具有重要的物理意义：观测中常见的慢传播对流耦合赤道扰动并不需要引入全新的动力学模态，而可以被理解为干波本征模态在热力加热作用下的连续变形。
CCEWs 与干波共享相同的模态结构与色散关系形式，差异仅在于有效等效深度的数值。
这一理论图景为 Wheeler 与 Kiladis 在波数-频率谱分析中观测到的结果提供了清晰的动力学解释：热带对流活动的显著谱峰精确地沿着 Matsuno 理论导出的色散曲线分布，只是等效深度取湿过程修正后的较小值。

为更系统地分析湿过程对赤道波动的影响，本文基于微扰理论进行了更进一步的研究。
在算子代数的表述中，非绝热加热被自然地表示为作用在干大气算子 $\Hop_0$ 上的微扰算子 $\Vop$，从而将复杂的热力过程转化为频率修正与模态混合等可计算的代数问题。

一阶微扰结果表明，频率修正等于微扰算子在未扰动态上的期望值。
对于 Wave-CISK 参数化，该修正使所有模态的频率按相同比例缩小，与精确解完全一致。
当加热与波动场之间存在相位差时，频率修正获得虚部，对应模态的增长或衰减率。
因此，波动的稳定性由加热算子的相位结构与本征态空间分布共同决定。

微扰理论还揭示了模态混合的选择定则。
加热微扰的空间对称性决定了其耦合不同经向模态的方式：关于赤道对称的加热仅耦合经向量子数相差为偶数的模态，而不对称加热则可耦合不同奇偶性的模态。
这种模态混合机制为解释观测中赤道波经向结构的复杂性提供了理论依据。

本文建立的理论框架虽然提供了理解赤道波动的统一视角，但仍存在若干局限性。
首先，将加热视为小参数微扰在强加热条件下可能失效，此时加热项与动力学项同阶，波导结构与本征模态可能发生根本性改变。
其次，本文中的加热被视为给定的参数化形式，忽略了加热与动力场之间的自洽反馈。
在强耦合极限下，这种反馈可能使系统行为超出线性分析的适用范围。
此外，垂直模态分解与等效深度假设在强对流条件下可能失效，不同垂直模态之间的耦合可能显著增强。

基于上述认识，未来研究可向多个方向拓展。
引入自洽的对流反馈参数化，系统研究从弱加热到强加热过渡过程中模态谱与稳定性的连续演化，是一个重要的理论方向。
放松线性假设，考察有限振幅条件下不同赤道模态之间的非线性相互作用及其对能量分布的影响，也具有重要意义。
将平均流、垂直切变及纬向非对称结构统一纳入算符化谱论框架，有望更直接地连接理想化理论与真实热带大气中复杂的对流-环流耦合过程。
这些方向的探索将进一步深化我们对热带大气动力学的理解，并为改进数值模式中的热带波动模拟提供理论指导。
